% frontmatter/como-usar.tex -- instrucciones para el adulto (en español).
\section*{Cómo usar este cuaderno}

Este cuaderno está pensado para una niña o un niño que ya reconoce las
letras y lee frases completas, y quiere ganar soltura practicando un
poco cada día -- no enseña a leer desde cero.

Tiene una página para cada día laborable del año: de lunes a viernes,
también en las vacaciones de Navidad y de verano -- es precisamente
cuando más se agradece no perder el ritmo de lectura. Cada página sigue
siempre la misma estructura.

\begin{itemize}[leftmargin=6mm]
\item \textbf{Fecha} y \textbf{Notas}: espacio para que un adulto anote
  la fecha y, si quiere, algo que observó ese día -- una palabra que
  costó, una que salió genial, el humor con el que llegó.
\item La \textbf{frase (o frases) del día}, dentro del recuadro verde.
  Léela en voz alta, despacio, señalando cada palabra con el dedo. Al
  principio del curso hay una sola frase; al final, cuatro. El texto se
  alarga cada trimestre, no cada semana -- no hay prisa.
\item Una \textbf{actividad} distinta cada día, siempre relacionada con
  lo que se acaba de leer:
  \begin{itemize}[leftmargin=6mm]
  \item \textbf{Dibuja}: dibuja algo de lo que cuenta la frase.
  \item \textbf{Completa}: hay unas pocas líneas ya dibujadas, sin forma
    todavía -- el resto lo inventa ella. No hay un dibujo ``correcto''.
  \item \textbf{Copia}: escribe la frase de hoy en las líneas, tres
    veces. Es para practicar la letra, no para inventar una respuesta --
    eso llega más adelante, cuando ya escribe sola (a partir del segundo
    trimestre verá también \textbf{Responde}: una pregunta corta sobre lo
    leído, para comprobar que se ha entendido y no solo que se ha leído).
  \item \textbf{Relaciona}: une con una línea las palabras que van
    juntas.
  \item \textbf{Adivina}: una adivinanza sencilla sobre la historia de
    la semana.
  \item \textbf{Crea}: una tarea abierta -- inventa, dibuja tu propia
    versión, cuenta lo que tú harías.
  \item \textbf{Repasa}: cada dos semanas, un pequeño resumen con
    casillas para marcar y una tarea final.
  \item \textbf{Traza}: de vez en cuando (más o menos cada dos
    semanas), en vez de dibujar, se repasa con el dedo o un lápiz el
    contorno de una letra -- mayúscula y minúscula -- y luego se
    escribe sola en la línea. Es para el trazo, no para la lectura: no
    hace falta ningún orden ni completarlas todas.
  \end{itemize}
\end{itemize}

Las páginas cuentan la historia de \textbf{Lucía}, su hermano
\textbf{Dani}, su perro \textbf{Toby} y el resto de su familia. Los
personajes se repiten durante todo el año, así que cada semana es un
capítulo nuevo de la misma historia.

No hace falta terminar la actividad el mismo día si no da tiempo -- lo
importante es la lectura. Cinco o diez minutos al día son suficientes.

Esta progresión está pensada para un ritmo ``medio''. Si durante una
semana entera lee la página sin ningún fallo a la primera, no hace
falta esperar al trimestre siguiente para ir más deprisa -- ese día se
puede leer más de una página. Si se atasca en casi cada palabra, que
lea solo la primera frase y que el resto lo lea un adulto en voz alta:
lo importante es que la sesión siga siendo agradable, no una lucha.

Al final de cada trimestre hay una pequeña medalla, no solo el diploma
del último día -- una meta más cercana que 260 días. Al final del
cuaderno hay también una \textbf{clave de respuestas} para las páginas
de Adivina y Relaciona.
\newpage
